\documentclass[11pt]{amsart}
\usepackage{amsmath,amssymb,amsthm}
\usepackage[margin=1.1in]{geometry}
\usepackage{booktabs}

\newtheorem{theorem}{Theorem}[section]
\newtheorem{lemma}[theorem]{Lemma}
\newtheorem{proposition}[theorem]{Proposition}
\newtheorem{corollary}[theorem]{Corollary}
\theoremstyle{definition}
\newtheorem{definition}[theorem]{Definition}
\newtheorem{remark}[theorem]{Remark}
\newtheorem{assumption}[theorem]{Computational Assumption}

\newcommand{\nn}{\mathbb{N}}
\newcommand{\vtwo}[1]{\nu_2(#1)}
\DeclareMathOperator{\HALT}{\textsf{HALT}}

\title{The Halting Problem of the Machine $1RB1LA\_0RC0RD\_1LC1LA\_1RE1RD\_1LF0LA\_{---}0LE$:\\ \large A Conditional Non-Halting Theorem and the Record-Timetable Obstruction}
\author{}
\date{October 2026}

\begin{document}
\maketitle

\begin{abstract}
We give a complete structural analysis of the $6$-state, $2$-symbol
Turing machine $M=$ \texttt{1RB1LA\_0RC0RD\_1LC1LA\_1RE1RD\_1LF0LA\_---\_0LE}
and prove, conditional on a single precisely-stated computational
assumption (master closure, Assumption~M), that it never halts from the
blank tape; unconditionally, non-halting holds over all runs whose round
classes stay in the exhaustively verified range ($h\le2000$, sampled to
$1.2\cdot10^{7}$) --- a doubly-exponential horizon whose step-count
scaling is labeled as empirical (Proposition on the envelope). The
analysis proceeds by
a macroscopic phase calculus between visits of state $D$ to blank cells,
a closed family of recurrent anchors, a halving-chain prologue, and a
\emph{round law} proved by strong induction on the $2$-adic class
$h=\nu_2(\alpha+1)$ of the leading counter, via a self-similar round
decomposition (halving cascade, universal ladder, master segment).
Magnitude decisions are handled unconditionally by a dominance inequality
(stated as Lemma~D of the companion derivation paper); the halting branch
is structurally unreachable for $h\ge82$
and perturbation-robust; $\Delta t(h)\in[2,9]$ on all verified classes.
The
single remaining assumption is shown to be equivalent, by a
finite-state reduction, to the continuation of a record timetable of the
autonomous $2$-adic walk, whose $32$-bit slice is certified acyclic by a
finite experiment (with non-eventual-periodicity beyond the slice stated
as a conjecture) --- the
precise, irreducible core of the problem. Two further groups of results,
unconditional
with respect to Assumption~M, are proved along the way: a mod-3 ledger --- from
step $54$ on, every phase configuration of $M$ has run-sum divisible by $3$, so
an entire half-space of configurations (including the Conway-chain entry
$S(3,1,3)$) is unconditionally unreachable and the re-entry and anchor
congruences hold as theorems --- and a complete birth-sequence theory:
a halt equivalence ($M$ halts iff some S-even birth value $A=x+y/2+1$ is a
power of two; both directions algebraic), sharpened by \emph{root reduction}
to ``iff some \emph{root} birth is a power of two'' (descendant births halve
exactly, preserving odd parts); \emph{all E1-reentry roots are killed
unconditionally} ($3\mid A$, a corollary of the ledger; $90.5\%$ of all
roots); a \emph{parity kill of dangerous keys} ($q\equiv1\bmod144\wedge q>1
\Rightarrow q\ge145\Rightarrow$ not a power; $19$ of the $20$ observed
classes permanently safe by key definition); and the \emph{drift law}
$\mathrm{M''}$ --- eight window families with the unified form
$A=3\cdot2^{k-1}w+2^{j}\mu$ at window positions
$13/26/50/99/110/194/219/220$, every family parity-killed, with the
$m\le40$ synthetic assertion table ($60$ segments, $0$ powers, $56/56$ exact
$F_{26}$ identities) --- which compresses the residual of Assumption~M into a
single arithmetic sentence in birth coordinates: the conditional chain of
the main theorem is thereby refined to ``never halts $\Leftarrow$
$\mathrm{M''}$ $+$ class preservation'', with the three precise gaps to an
unconditional proof (window-enumeration completeness $=$ M; conditionality of
class preservation; $m>25$ as a constructed-anchor property) stated
explicitly. Per-depth raw-reproduced witnesses (through depth $3$) show
that birth guards fail at every finite depth, stated at witnessed level. All
other components are proved
algebraically or verified exhaustively over declared finite domains;
every claim is classified in a ledger (algebraic / exhaustive /
assumption-dependent).
\end{abstract}

\section{Introduction}\label{s:intro}

Let $M$ denote the Turing machine with states $A,\dots,F$ and transition
table of Table~\ref{tab:tm}, in the standard notation where, e.g.,
\texttt{1RB} means: write $1$, move right, go to state $B$, and where the
single undefined transition (state $F$ reading $0$) is the halting
transition. The machine is one of the remaining undecided instances of the
$6$-state Busy Beaver problem; to our knowledge no analysis of it has been
published.

\begin{table}[h]
\centering
\begin{tabular}{c|cc}
 & read $0$ & read $1$\\\toprule
$A$ & \texttt{1RB} & \texttt{1LA}\\
$B$ & \texttt{0RC} & \texttt{0RD}\\
$C$ & \texttt{1LC} & \texttt{1LA}\\
$D$ & \texttt{1RE} & \texttt{1RD}\\
$E$ & \texttt{1LF} & \texttt{0LA}\\
$F$ & $\HALT$ & \texttt{0LE}\\
\end{tabular}
\caption{The transition table of $M$.}
\label{tab:tm}
\end{table}

\begin{theorem}[Main, conditional form]\label{thm:main}
$M$, started on the blank tape, never halts, conditional on
Assumption~\ref{ass:master} (master closure for all classes).
\end{theorem}

A second form of the conditional chain --- ``never halts $\Leftarrow$ drift
law $\mathrm{M''}$ $+$ class preservation'' --- is established in
\S\ref{s:ledger} (Corollary~\ref{cor:assembly2}); by
Proposition~\ref{prop:drift}(d) the two forms carry the same non-halting
content within the established architecture.

\begin{proposition}[unconditional envelope, verified range; empirical status made explicit]\label{thm:envelope}
Unconditionally, $M$ does not halt within any run whose round classes stay
in the exhaustively verified range (all classes $h\le2000$; sampled
agreement through $1.2\cdot10^{7}$) --- in step-count terms, non-halting
is proven over a horizon that is doubly exponential in the class index
(below $2^{2^{2000}}$ in the record-climb heuristic).
\end{proposition}
\begin{proof}[Status of the envelope]
The finite part is unconditional: over the verified range the round table
is single-path and carries no halting block (Theorem~\ref{thm:round},
computational components of \S\ref{s:comp}); the anchors climb by
$n'\ge n+6$ per exit, and the record timetable of the residual climbs by
a measured factor $\approx4$ per class level, so the first arrival of
class $h$ requires a number of rounds growing like $4^{h}$, each of
supercubic cost in the anchor index --- hence the doubly-exponential
horizon. The record climb, however, is a \emph{measured} law of the
autonomous walk (Theorem~\ref{thm:noCycle}), not an unconditional
theorem; we therefore state the envelope as a proposition whose
step-count translation carries this explicitly labeled empirical
component. A paper proof of an explicit lower bound
$T(h)\ge2^{c\cdot2^{h}}$ (assembling $A_j\le\tau_j^2$ with the band
$\Delta\alpha\le2h+6$) is left open. Independently of the class range, the
shape of the unconditional envelope is sharpened by
Theorems~\ref{thm:mod3} and~\ref{thm:haltpow}: every configuration with
$\Sigma\not\equiv0\pmod3$ is unreachable from step $54$ on (an entire
half-space, with no appeal to any verified class range), and ``$M$ never
halts'' is equivalent to ``no power-of-two S-even birth ever occurs'' --- the
verified domain of the latter equivalence being the first $10^{7}$ phase
events ($2{,}908$ births, none a power of two; grade B); by root reduction
(Theorem~\ref{thm:root}) it is further equivalent to ``no power-of-two
\emph{root} birth ever occurs'', whose verified domain is: $30{,}791{,}021$
root births over $10^{7}$ anchors with $0$ powers, all $20$ observed classes
safe ($18$ by key definition $+$ Lk$|$j3 single-point $+$ L3$|$j4 family
form), all eight dangerous window families killed through constructed
anchors $a\le300$, and the $m\le40$ assertion table all passing (grade B).
\end{proof}

The proof occupies Sections~\ref{s:phase}--\ref{s:assembly}. Its skeleton:
\begin{enumerate}
\item (\S\ref{s:phase}) A phase calculus for the instants when the head, in
state $D$, reads a $0$. Between two such instants the machine executes a
fully determined finite block of elementary steps, and the possible tape
shapes at these instants form a finite grammar with ten transition rules
(nine non-halting event edges and the halting block; Lemma~\ref{lem:rules}).
The unique halting path passes through a single
grammar configuration $G(3,\cdot)$ (Lemma~\ref{lem:halt}).
\item (\S\ref{s:prologue}) The blank-tape run reaches at step $54$ an
\emph{anchor} $P(3)$, and every anchor $P(n)$ is carried, by an explicit
halving chain of length $m=\nu_2(n+3)$, into a standard family
$\Phi(w,m)$, $n+3=2^m w$, $w$ odd (Theorem~\ref{thm:prologue}).
\item (\S\ref{s:ledger}) Unconditionally with respect to
Assumption~\ref{ass:master}: the mod-3 ledger (an entire half-space of
configurations unreachable from step $54$ on), the halt equivalence via
power-of-two births and its root-reduction sharpening, the E1-reentry kill,
the parity kill of dangerous keys, the L-exit class snapshot, and the drift
law $\mathrm{M''}$ --- yielding the second form of the conditional chain,
``never halts $\Leftarrow$ $\mathrm{M''}$ $+$ class preservation''
(Corollary~\ref{cor:assembly2}), with the three precise gaps stated
(Remark~\ref{rem:gaps}).
\item (\S\ref{s:landmark}--\ref{s:ladder}) The $\Phi$-trajectory visits
\emph{landmarks} $L(\alpha,t;s)$ indexed by a leading counter
$\alpha\equiv1\pmod 6$ and a time parameter $t$; the evolution of
$(\alpha,t)$ depends only on the class $h=\vtwo{\alpha+1}$
(Theorem~\ref{thm:round}, the \emph{round law}), proved by strong induction
on $h$ via a halving cascade (Lemma~\ref{lem:cascade}), a universal ladder
(Lemma~\ref{lem:ladder}), and a uniform master segment
(Assumption~\ref{ass:master}, uniformity clause --- Remark~\ref{rem:uniform}).
\item (\S\ref{s:endgame}--\ref{s:assembly}) Exhausted-tail configurations
exit to the next anchor with $n'\ge n+6$ and $n'\equiv3\pmod 6$
(Theorem~\ref{thm:endgame}), and no reachable configuration is the halting
one (Theorem~\ref{thm:safety}). The anchors therefore climb forever.
\end{enumerate}

Throughout, tape contents are described by \emph{runs} (maximal blocks of
$1$s) separated by \emph{gaps} (blocks of $0$s); all configurations live on
the right-infinite half of the tape beyond a fixed left sentinel, which the
machine never crosses (this follows from the left-moving transitions
returning before the sentinel; it is confirmed by every invariant below).

\section{The phase calculus}\label{s:phase}

\subsection{Events and configurations}

\begin{definition}[event]
An \emph{event} is an instant at which the head, in state $D$, reads a $0$.
Let $\mathcal{C}$ denote the set of tape shapes (with head position) that
can hold at an event.
\end{definition}

\begin{definition}[grammar]\label{def:grammar}
Write a configuration as a list of runs $(r_1,r_2,\dots,r_k)$,
$k\ge2$, with specified gaps between them and a marked head cell (the $0$
being read). The following four types exhaust the configurations that arise
on the orbits considered:
\begin{itemize}
\item $R(a,b)$: runs $(a,b)$, gaps $[1]$, head on the fresh $0$
\emph{immediately right} of the right run;
\item $S_k(x,y,z,\dots)$: $k\ge3$ runs, gaps $[1,2,1,2,\dots]$, head on the
\emph{left} $0$ of the $2$-gap between runs $2$ and $3$;
\item $G_k(a,b,\dots)$: $k\ge2$ runs, gaps $[3,1,1,\dots]$, head on the
\emph{middle} $0$ of the $3$-gap between runs $1$ and $2$;
\item $L_k(r_1,r_2,r_3,\dots)$: $k\ge3$ runs, gaps $[1,1,1,\dots]$, head on
the single-$0$ gap between runs $2$ and $3$.
\end{itemize}
(Thus $L_2$ is identified with $R$: both have the head just right of the
block.)
\end{definition}

\begin{lemma}[phase rules]\label{lem:rules}
Between consecutive events, $M$ executes exactly the following transitions,
with the stated exact step counts:
\[
\begin{array}{llc}
\text{(R)} & R(a,b)\to S_3(a{+}1,\,b{-}2,\,1) & 2b+3,\quad b\ge3;\\[2pt]
\text{(S$\ge$)} & S_k(\dots,x,y,z,\dots)\to S_k(\dots,x{+}1,\,y{-}2,\,z{+}1,\dots) & 2y+3,\quad y\ge3;\\[2pt]
\text{(S2)} & S_k(\dots,x,2,z,\dots)\to G_{k-1}(\dots,x{+}1,\,z{+}1,\dots) & 7;\\[2pt]
\text{(S1)} & S_k(\dots,x,1,z,\dots)\to L_{k-1}(\dots,1,\,x{+}z{+}3,\dots) & 2x+z+14;\\[2pt]
\text{(G$\ge$)} & G_k(a,b,\dots)\to S_{k+2}(1,\,a{-}3,\,1,\,b{+}1,\dots) & 2a+3,\quad a\ge4;\\[2pt]
\text{(G3)} & G_k(3,b,\dots)\to G_{k+1}(1,\,1,\,b{+}1,\dots) & 9;\\[2pt]
\text{(G2)} & G_k(2,b,\dots)\to L_{k+1}(1,\,3,\,b{+}1,\dots) & 16;\\[2pt]
\text{(G1)} & G_k(1,b,\dots)\to \HALT & 6;\\[2pt]
\text{(L$\ge$)} & L_k(r_1,r_2,r_3,\dots)\to L_k(r_1{+}1,\,r_2,\,r_3{-}1,\dots) & 2r_2+5,\quad r_3\ge2;\\[2pt]
\text{(L1)} & L_k(r_1,r_2,1,\dots)\to R/S_{k-1}(r_1{+}1,\,r_2,\dots) & 2r_2+5.
\end{array}
\]
In (S1) the result is $R$ when $k=3$ and $L_{k-1}$ when $k\ge4$; in (L1) it
is $R$ when $k=3$ and $S_{k-1}$ when $k\ge4$. The halting block (G1) ends
with $F$ reading $0$; the block (G3) is an ordinary event block whose target
$G(1,\cdot)$ halts in the next block.
\end{lemma}

\begin{proof}[Proof sketch]
Each rule is obtained by executing the transition table from the stated
configuration until the next event. We illustrate rule (R) for $R(3,3)$
(tape $\dots1^3\,0\,1^3\,0^\infty$, head on the fresh $0$, state $D$;
the trace below is a step-level recording of the machine):
\[
\begin{array}{c|ll}
t & \text{state/head} & \text{effect}\\\hline
0 & D@7 & \text{tentative write }1\text{, right, }E\\
1 & E@8 & \text{write }1\text{, left, }F\\
2 & F@7 & \text{\emph{erase} (}0\texttt{LE), left, }E\\
3 & E@6 & \text{\emph{erase} (}0\texttt{LA), left, }A\\
4,5 & A@5,A@4 & \text{march left through the right run}\\
6 & A@3 & \text{read the gap }0\text{: \emph{fill} it (}1\texttt{RB), right, }B\\
7 & B@4 & \text{write }0\text{ (}0\texttt{RD), right, }D\\
8 & D@5 & \text{march right writing }1\text{s}\\
9 & D@6 & \text{next event: }S_3(4,1,1)
\end{array}
\]
The net effect is a five-phase mechanism: (i) $D$'s tentative extension of
the right run is \emph{undone} by the $E$--$F$--$E$--$A$ recoil; (ii) $A$
marches left across the right run; (iii) $A$ fills the inter-run gap with
$1$, growing the first run to length $a{+}1$; (iv) $B$ zeroes one cell of
the old right run, creating the new $2$-gap; (v) $D$ marches right,
rebuilding, so that the old right run of length $b$ leaves a residue $b-2$
to the left of the new gap and a fresh run of length $1$ to its right:
$S_3(a+1,b-2,1)$, at total cost $2b+3$. The remaining rules follow
by the same kind of shuttling; each was derived stepwise from
Table~\ref{tab:tm}. Rule-level faithfulness is machine-checked in two
complementary ways: (i) the rule checker \texttt{rulecheck.mjs} compares
every rule against the raw machine on parameter grids ($a\in\{1,\dots,20\}$,
$b\in\{2,3,5,10\}$, tails of length $0$--$2$): $564/564$ rows agree exactly
(this includes the boundary cases $y\in\{1,2\}$, $a\in\{1,2,3\}$, and the
halting case (G1)); (ii) trajectory-level comparison of the phase engine
against the raw oracle over $956$ anchor trajectories ($n=2..957$):
$955$ agree field-by-field, including $16$ halting anchors with exact halt
times; the single exception is $R(2,2)$, whose successor ($b=2<3$) lies
outside the declared grammar and is flagged, not mis-computed (direct
simulation halts at $+22$). An earlier version of the table carried a
folded (G3)$\to\HALT$ row and an incorrect (G2) image
$L(1,b{+}1,b{+}1)$; against that version the checker reported
$531$ ok / $21$ bad, the bad rows being precisely (G2)/(G3) --- the
correction here follows the raw-machine side of that checker.
\end{proof}

\begin{lemma}[halting criterion]\label{lem:halt}
$M$ halts from an event configuration iff the phase trajectory reaches
$G_k(1,b,\dots)$ for some $b$ and tail; equivalently, iff it reaches a
$G$-state with leading head $a\in\{1,3\}$ (rule (G3) carries $G(3,\cdot)$
to $G(1,\cdot)$ in one block). On every orbit considered below, the
leading heads of all $G$-states are $\ge2$ and $\ne3$, so no other rule
can produce a halt.
\end{lemma}

\begin{proof}
The only undefined transition of Table~\ref{tab:tm} is $F$ reading $0$. An
exhaustive check of the ten rules shows the state $F$ reads a $0$ in the
interior of an event block in exactly one grammar situation, namely the
configuration $G(1,b,\dots)$ (rule (G1)); in all other blocks $F$ reads
$1$s or the block ends at the next event first. Rule (G3) is the sole
in-rule source of $G(1,\cdot)$.
\end{proof}

\subsection{Sweeps}

For efficiency we also use \emph{swept} versions of the rules: consecutive
applications of (S$\ge$) with the same $x$-growth collapse into one step
with closed-form cost. If $y$ is odd, $j=(y-1)/2$ applications carry
$(x,y,z)\mapsto(x+j,\,1,\,z+j)$ at total cost $\sum_{i=0}^{j-1}(2(y-2i)+3)$;
if $y$ is even, $j=(y-2)/2$ applications carry $(x,y,z)\mapsto(x+j,\,2,\,z+j)$.
These identities are elementary and are used silently below.

\section{Anchors and the prologue}\label{s:prologue}

\begin{definition}[anchor]\label{def:anchor}
For $n\ge2$ let $P(n)$ denote the configuration
$0^\infty\,1^n\,0\,1^n\,0^\infty$ with the head in state $D$ on the fresh
$0$ immediately right of the right run; in grammar terms $P(n)=R(n,n)$.
\end{definition}

\begin{lemma}[entry]\label{lem:entry}
Started on the blank tape, $M$ reaches $P(3)$ at step $54$.
\end{lemma}

\begin{proof}
Direct simulation of the $54$ steps: the machine builds the pattern
$1^3\,0\,1^3$ by the $A$-$B$-$C$-$D$-$E$-$F$ loop and lands in state $D$ on
the fresh $0$. (Step-level trace in the repository; verified by two
independent simulators.)
\end{proof}

\begin{definition}[the family $\Phi$]\label{def:phi}
For odd $w\ge3$ and $m\ge1$:
\[
\Phi(w,m)\;=\;G\bigl(w-1,\;w-1,\;2w,\;4w,\;\dots,\;2^{m-1}w,\;2^{m}w-1\bigr).
\]
\end{definition}

\begin{theorem}[prologue]\label{thm:prologue}
Let $n\equiv3\pmod6$, $n\ge3$ (this holds for every anchor on the orbit of
Theorem~\ref{thm:main}; in particular $n$ is odd), and write
$n+3=2^{m}w$ with $w$ odd, $m\ge1$. Then the phase
trajectory from $P(n)$ reaches $\Phi(w,m)$; explicitly,
\[
R(n,n)\;\to\;R(1,\,2n+2)\;\to\;G(n+2,\,n+1)\;\to^{\,m\ \text{halvings}}\;\Phi(w,m),
\]
where the halving steps act on the two leading runs while the tail is
carried through, and the leading heads follow the closed form
\begin{equation}\label{eq:halving}
a_i \;=\; 2^{m-i}\,w-1,\qquad i=0,1,\dots,m .
\end{equation}
\end{theorem}

\begin{proof}
Since $n$ is odd, rule (R) gives $R(n,n)\to S_3(n+1,\,n-2,\,1)$ with
middle run $y=n-2$ \emph{odd}; the sweep drives $y$ down to $1$ in
$j=(n-3)/2$ steps, and rule (S1) yields $R(1,\,x+z+3)$ with
$x=n+1+j$, $z=1+j$, i.e.\
\[
R(1,\,2n+2).
\]
Rule (R) on $b=2n+2$ gives $S_3(2,\,2n,\,1)$ with $y=2n$ \emph{even}; the
sweep drives $y$ down to $2$ in $n-1$ steps, and rule (S2) yields
\[
G\bigl(2+(n-1)+1,\;1+(n-1)+1\bigr)\;=\;G(n+2,\,n+1).
\]

Now let $G(a,b,\mathcal R)$ be a halving state of the chain started at
$G(n+2,n+1)$, i.e.\ $a+1=2^{m-i}w$ with $0\le i\le m$. Since $w$ is odd,
$a$ is odd exactly when $i<m$, so the halving map (rules (G$\ge$) followed
by the $L/S$ shuttles of \S\ref{s:phase}, whose composite on the two
leading runs is
$(a,b)\mapsto\bigl((a-1)/2,\,(a-1)/2\bigr)$ with $b$ entering the
accumulating geometric tail) applies for $i=0,\dots,m-1$. Starting from
$a_0+1=n+3=2^{m}w$ the heads satisfy \eqref{eq:halving}, terminating at
$a_m=w-1$ with accumulated tail $(2w,4w,\dots,2^{m-1}w)$ and terminator
$2^{m}w-1$: exactly $\Phi(w,m)$.
\end{proof}

\begin{remark}
Theorem~\ref{thm:prologue} is stated for $n\equiv3\pmod6$, the congruence
class of every anchor reached from $P(3)$ (Theorem~\ref{thm:endgame},
\S\ref{s:assembly}); for even $n$ the initial sweep branches differently
and the chain above does not apply. The theorem is algebraic and holds for
all $m\ge1$; it was additionally machine-witnessed for all
$n\equiv3\pmod6$, $n\le6003$ ($1001$ values), and for
$m\in\{13,\dots,1200\}$, $w\in\{3,9,21\}$.
\end{remark}

\section{The mod-3 ledger and the birth sequence}\label{s:ledger}

This section contributes a chain of results that are \emph{unconditional} with respect to
Assumption~\ref{ass:master}: a mod-3 ledger for the whole orbit
(Theorem~\ref{thm:mod3}), a halt equivalence in terms of power-of-two births
(Theorem~\ref{thm:haltpow}) with its root-reduction sharpening
(Theorem~\ref{thm:root}), unconditional safety of all E1-reentry root births
(Lemma~\ref{lem:e1}), a per-key parity kill (Theorem~\ref{thm:parity}), the
L-exit class-table snapshot (Proposition~\ref{prop:classes}), the m-record
timetable (Proposition~\ref{prop:mrecords}), the drift law $\mathrm{M''}$ ---
the final explicit form of the residual (Proposition~\ref{prop:drift}) --- and
the synthetic assertion table through $m=40$
(Proposition~\ref{prop:m2640}); witnessed-level propositions on
the pollution of birth guards follow
(Propositions~\ref{prop:depth1} and~\ref{prop:cascade}). Together they seal the
mod-3 side unconditionally, kill all E1-reentry roots, kill every identified
L-exit family by pure parity, and condense the entire residual content of
halt-exclusion into the single statement $\mathrm{M''}$, whose precise
three-ring gap to an unconditional proof is spelled out
(Remark~\ref{rem:gaps}).

\begin{definition}[ledger invariant $I_0$]\label{def:I0}
For a configuration $C$ let $\Sigma(C)$ denote the sum of its runs. $I_0$ is
the set of configurations satisfying: (i) $\Sigma\equiv0\pmod3$; (ii) domain:
$R(a,b)$ with $b\ge3$; $S(x,y,z,\overline r)$ with $y\ge1$; $G(a,b,\overline
r)$ with $a\ge2$; $L(r_1,r_2,r_3,\overline r)$ with $r_2\ge3$ and $r_3\ge1$;
(iii) boundary guards BG1: $G(2,b,\overline r)$ requires $b\ge2$ and
$b\equiv2\pmod3$; BG2: $S(1,2,z,\overline r)$ requires $z\equiv1\pmod3$.
Both guards are \emph{forced} by closure: BG2 because (S2) births
$G(2,z{+}1)$, which must pass BG1; BG1 because the $\Sigma$-increment of the
(G2) edge in the derivation engine's form is $b+1$ (at the raw level the
increment is constantly $+3$; Remark~\ref{rem:g2}). Under
$\Sigma\wedge\mathrm{BG2}$ the three-run class $S(1,2,z)$ is empty
($z\equiv0$ contradicts $z\equiv1$); the class $S(1,2,z,\overline r)$ with
$k\ge4$ is nonempty by tail freedom, and the orbit runs precisely there
((G$\ge$) at $a=5$ births $S(1,2,1,b{+}1,\overline r)$, $z=1$).
\end{definition}

\begin{lemma}[unique reachable value of $G(2,\cdot)$]\label{lem:g22}
The grammar-reachable configurations $G(2,b,\overline r)$ are exactly those
with $b=2$ (including the geometric-tail family $\Phi(3,m)$).
\end{lemma}
\begin{proof}
The only in-rule source of a $G$-state is (S2),
$S(x,2,z,\overline r)\to G(x{+}1,z{+}1,\overline r)$; head $2$ forces $x=1$ at
the sweep terminal. An $S$-entry with $x=1$ can only come from (G$\ge$)
((R) births $x=a{+}1\ge2$; (L1) births $x=r_1{+}1\ge2$; (S$\ge$) sweeps
increase $x$), i.e.\ $S(1,a{-}3,1,b{+}1,\overline r)$ with $z=1$ fixed;
reaching $y=2$ with $x$ still $1$ forces the sweep length $j=0$, i.e.\ $a=5$,
yielding $S(1,2,1,b{+}1)\to G(2,2,b{+}1,\overline r)$.
\end{proof}

\begin{theorem}[mod-3 ledger; unconditional w.r.t.\ Assumption~M]\label{thm:mod3}
(a) \emph{Closure.} Every member of $I_0$ steps to a member of $I_0$ under
each of the ten rules of Lemma~\ref{lem:rules}, the only exceptions being the
halt forms $G(3,b,\overline r)$ (classified, not violations).
(b) \emph{Seed.} The blank-tape orbit enters $R(3,3)\in I_0$ at step $54$
($\Sigma=6$). (c) \emph{Conclusion.} Every phase (grammar) configuration of
$M$ from step $54$ on satisfies $\Sigma\equiv0\pmod3$. Consequently every
re-entry $R(1,b)$ has $b\equiv2\pmod3$; every anchor $R(n,n)$ has
$n\equiv0\pmod3$; any halt must occur at a $\Sigma\equiv0$ form
$G(3,b,\overline r)$; and every configuration with $\Sigma\not\equiv0$ ---
including the Conway-chain entry $S(3,1,3)$ ($\Sigma=7$) --- is unreachable,
an entire half-space. $I_0$ does not exclude halting (the halt forms
$G(3,\cdot)$ can lie in $I_0$); the entire residual content of
halt-exclusion lies on the 2-adic birth side (Theorem~\ref{thm:haltpow}).
\end{theorem}
\footnote{Granularity of the halt segment: the folded block
(G3)$+$(G1) ($15$ steps total) contains, at $+9$ steps, the intermediate
event $G(1,1,b{+}1,\overline r)$ --- the target of (G3), outside the
$G$-domain ($a\ge2$) of $I_0$ but still with $\Sigma\equiv0$ (the (G3)
$\Sigma$-identity has $\Delta=0$; spot-verified against the raw machine). The
universal statement of Theorem~\ref{thm:mod3}(c) is thus pinned to
\emph{phase (grammar) configurations}.}

Proof status (grade A$^-$, closure part): the closure is given by per-edge
elementary affine identities plus guard transfers --- (R):
$(a{+}1)+(b{-}2)+1=a+b$; (S$\ge$) and (L$\ge$): translations; (S2):
$(x{+}1)+(z{+}1)=x{+}z{+}2$; (S1), both branches: $+3$; (G$\ge$):
$1+(a{-}3)+1+(b{+}1)=a+b$; (G3): $\Delta=0$; (G2): Remark~\ref{rem:g2};
(L1), both branches: $+1$ --- each hand-checkable, and re-verified by machine
sweeps: the author's boundary enumeration of $219{,}398$ $I_0$ member-steps
plus $2^{40}$-scale random member-steps (the assertion total fluctuates by
about $\pm60$ with the random sampling; the deterministic core reproduces
digit-for-digit), and an independent heterogeneous closure grid of $20{,}508$
member-steps with run counts up to $k\le10$ and $a=0$ probes. The reduction
of closure to bounded run count is carried by locality (each rule rewrites
only the first three runs; both guards read only the first three).
Event-level raw fidelity was independently re-verified (all ten rules,
boundary grids plus $2^{24}$-scale random); the seed is triple-sourced (two
simulators plus an independent raw-machine parse). The orbit side is a
redundant confirmation (the induction already yields the all-times
statement): $10^{7}$ phase events with zero violations ($406$ anchors, all
$n\equiv0$; $1{,}485$ re-entries, all $b\equiv2$); an independent re-check of
$5\times10^{5}$ events with zero violations ($91$ anchors, $332$ re-entries);
the trajectory library of $1{,}357{,}014$ states with zero violations
($R$-family decomposition $825=547+150+128$, all $a+b\equiv0$). The $112$
halt forms $G(3,\cdot)$ mentioned above are an enumeration count of the
closure grid, not an orbit observation (the orbit count of $G(3,\cdot)$ is
$0$).

\begin{remark}[raw fidelity domain of (G2)]\label{rem:g2}
Raw truth: the (G2) edge maps $G(2,b,\overline r)\to L(1,3,b{+}1,\overline
r)$ for \emph{all} $b\ge2$ (raw-verified for $b\in\{2,3,5,8,11,14\}$ with
three tail shapes; verbatim the (G2) row of Lemma~\ref{lem:rules}, whose
rule checker $564/564$ also covers (G2) at each $b$), with $\Sigma$-increment
constantly $+3$ --- so the mod-3 ledger is preserved by the true machine on
all of $I_0$ \emph{without appeal to BG1} (under the raw identity BG1 becomes
superfluous; only the engine-level bookkeeping needs it). The derivation
engine of the ledger analysis suite carries the simplified image
$L(1,b{+}1,b{+}1,\overline r)$, which agrees with the raw machine precisely
at $b=2$; by Lemma~\ref{lem:g22} this is the unique grammar-reachable value,
so the distortion site $G(2,b\neq2)$ is unreachable, and even if it were
entered (started from an off-grammar configuration) the increment $+3$ would
still preserve the ledger. Accordingly, the reading ``ten rules $=$ raw
machine'' must be qualified: the rule table is event-level raw-faithful on
the grammar domain; on the off-grammar layer ($G(2,b\neq2)$) the engine uses
the simplified image, with no effect on the ledger.
\end{remark}

\begin{definition}[S-even birth]\label{def:birth}
An event whose configuration is an $S(x,y,z,\overline r)$ with $y$ even
($y\ge2$) is called an \emph{S-even birth} (the entry of that even sweep);
its \emph{birth value} is
\[
A \;=\; x+\tfrac{y}{2}+1 .
\]
\end{definition}

\begin{theorem}[halt equivalence via power-of-two births]\label{thm:haltpow}
$M$ (started on the blank tape) halts if and only if some S-even birth value
$A=x+y/2+1$ on its phase trajectory is a power of two; equivalently
$A\in\{4,8,16,\dots\}$ (always $A\ge3$, since $x\ge1$ and $y\ge2$; for $A=3$
the sweep terminal has head $2$ and feeds the non-halting (G2) channel).
\end{theorem}
\begin{proof}
First, the \emph{sweep invariant}: (S$\ge$) preserves $A$, since
$(x{+}1)+\frac{y-2}{2}+1=x+\frac y2+1$; a birth sweep terminates at
$S(x',2,z')$ with $x'=x+\frac{y-2}{2}$, and (S2) lands in the $G$-state with
head $x'{+}1=A-1$.
($\Leftarrow$) Let some birth have $A=2^{t}$ ($t\ge2$). Its sweep terminal
$G$-state has head $2^{t}-1$ (odd, $\ge3$). For $t=2$ this is already the
halt form $G(3,\cdot)$. For $t\ge3$ the head is $\ge7$ and odd, and the
halving composite of Theorem~\ref{thm:prologue} (on the two leading runs,
$(a,b)\mapsto\bigl(\frac{a-1}{2},\frac{a-1}{2}\bigr)$ with the tail carried
verbatim) applies at each step: the head chain runs
$2^{t}-1\to2^{t-1}-1\to\cdots\to7\to3$ (every head on the chain has the form
$2^{s}-1\ge3$, never the value $5$, never the $G(2,\cdot)$ channel), ending
at the halt form $G(3,\cdot)$; by Lemma~\ref{lem:halt}, $M$ halts.
($\Rightarrow$) Suppose $M$ halts. By Lemma~\ref{lem:halt} the trajectory
reaches a $G$-state with head $3$; the only in-rule source of
$G(3,z{+}1,\overline r)$ is (S2) acting on $S(2,2,z,\overline r)$. That
$S$-state has $y=2$ and lies inside some even sweep (or is itself its entry);
the sweep invariant gives the birth value at that sweep's entry
$A=2+\frac22+1=4=2^{2}$, a power of two.
\end{proof}

Proof status (grade A$^-$, both directions): both directions are elementary
algebra, resting on rule faithfulness and the uniqueness of the halt form
(heads $\in\{1,3\}$ with only the (G3)/(G1) exits, raw-verified broadly);
rule faithfulness is the B-grade check (rulecheck $564/564$, $956$ anchors)
plus the independent event-level raw re-verification. Empirical companion
(grade B): in the first $10^{7}$ phase events there are $2{,}908$ S-even
births, none a power of two; an independent re-check of $5\times10^{5}$
events finds $647$ births, none. The entire content of Assumption~M condenses
to the single sentence ``no power-of-two birth ever occurs''. The next two
subsections compress that sentence further, onto \emph{root} births and then
onto the dangerous keys of L-exits, ending with the explicit form
$\mathrm{M''}$ of the residual.

\subsection{Root reduction and E1-reentry safety}

\begin{definition}[root and descendant births]\label{def:rootbirth}
A birth (Definition~\ref{def:birth}) is exactly one of the following by its
entry edge: an \emph{E1 root}, entering the even sweep via (R) from $R(a,b)$
with $b$ even, $A=a+b/2+1$ --- on the orbit, the re-entries $R(1,b)$ ($b$ even)
and the anchor initializations $R(n,n)$ ($n$ even); an \emph{L-exit root},
entering via (L1) from $L(r_1,r_2,\overline r)$ with $r_2$ even,
$A=\sigma+r_2/2+1$ with $\sigma=r_1+r_3$ ($L_3$ relays through
$R(\sigma,r_2)$ with the same formula: $L_3\to R(\sigma,r_2)\to
S(\sigma{+}1,r_2{-}2,1)$, $A=\sigma+1+\frac{r_2-2}{2}+1=\sigma+r_2/2+1$); or a
\emph{descendant}, entering via (G$\ge$) from $G(a,\overline r)$ with $a$ odd
$\ge5$, of value $A'=(a{+}1)/2$ --- where $a{+}1$ is exactly the birth value
that produced that $G$-state (the (S2) landing head is $A-1$; sweep invariant
in the proof of Theorem~\ref{thm:haltpow}).
\end{definition}

\begin{theorem}[root reduction]\label{thm:root}
Every descendant birth has exactly half the value of its chain's root birth;
inductively, \emph{the odd part of every birth equals the odd part of its
chain's root birth}. In particular,
\[
M\ \text{halts}\iff\ \text{some \emph{root} birth has } A=2^{t}\ (t\ge2).
\]
\end{theorem}
\begin{proof}
By Definition~\ref{def:rootbirth} descendant values $A'=(a{+}1)/2=A/2$ halve
exactly, so odd parts are chain-invariant.
($\Leftarrow$) Let a root birth have $A=2^{t}$: its even sweep lands via (S2)
in the $G$-state with head $2^{t}-1$, after which the halving composite
(Lemma applied as in Theorem~\ref{thm:prologue}) runs
$2^{t}-1\to2^{t-1}-1\to\cdots\to7\to3$, reaching the halt form $G(3,\cdot)$
after $t-2$ steps --- every head on the chain has the form $2^{s}-1\ge3$,
independently of $b$ and the tail (carried verbatim); by Lemma~\ref{lem:halt}
$M$ halts.
($\Rightarrow$) This is the ($\Rightarrow$) of Theorem~\ref{thm:haltpow}: a
halt gives some S-even birth equal to $2^{t}$; if that birth is a descendant
its odd part is $1$, hence its chain root has odd part $1$, i.e.\ the root is
a power $2^{s}$.
\end{proof}

\begin{proof}[Proof status (grade A$^-$)]
The classification algebra is as above (resting on rule faithfulness and the
halving composite); the exhaustive machine evidence for the provenance
classification (grade B, adversarial re-audit ref-b1): over $2\times10^{7}$
uncollapsed events / $572$ anchors, all $4{,}120$ births originate from R/L/G
only ($8{,}576{,}839$ S-interior repeats correctly deduplicated), every
$R$ tagged $L_3$ comes from the $L_3\to R$ relay; every $G$-entry head
$a{+}1$ equals the S-entry birth value that produced it ($0$ violations);
descendant inheritance $A'=(a_{\mathrm{prev}}{+}1)/2$ ($0$ violations);
descendant odd part $=$ nearest root odd part ($0$ violations); the
uncollapsed and collapsed engines agree term-by-term on the $(\mathrm{ty},A)$
root sequences ($0$ mismatches; births10's $600$-anchor cross-check
independently reproduced); an independent raw-machine simulator reproduces
the first $12$ anchor transitions of the blank-tape run exactly in both
$n'$ and cost.
\end{proof}

\begin{lemma}[E1-reentry root safety]\label{lem:e1}
Every re-entry $R(1,b)$ on the orbit satisfies $b\equiv2\pmod3$
(Theorem~\ref{thm:mod3}, unconditional). If moreover $b$ is even (the only
case that births), then $b\equiv2\pmod6$, $b\ge8$, and
\[
A=\tfrac b2+2\ \equiv\ 1+2\equiv0\pmod3,\qquad A\ge6 .
\]
Hence \emph{no E1-reentry root birth is ever a power of two} ---
unconditionally with respect to Assumption~M.
\end{lemma}
\begin{proof}
$b\equiv2\pmod6$ gives $b/2\equiv1\pmod3$; the domain $b\ge3$ with
$b\equiv2\pmod6$ forces $b\ge8$, so $A\ge6$; and $2^{t}\not\equiv0\pmod3$.
\end{proof}
Empirical companion (grade B, independent counter ref-b2): over $10^{6}$
anchors, all $2{,}787{,}614$ E1-reentry root births satisfy $3\mid A$ ($0$
exceptions; all with $b\equiv2\pmod6$; least odd part $3$); all
$3{,}669{,}365$ re-entries (either parity of $b$) have $b\equiv2\pmod3$.

\begin{lemma}[sharpened AP power criterion]\label{lem:ap}
Let $D=2^{s}u$ ($u$ odd) $>0$ and $A_0=2^{j}q_0$ ($q_0$ odd) with $j<s$.
Then the arithmetic progression $\{A_0+Dk:k\ge0\}$ contains a power of two
\emph{iff} $q_0=1$ (i.e.\ $A_0$ itself is a power). For $j\ge s$ containment
is decidable: $\exists\,t\ge s$ with $2^{t}\ge A_0$ and $2^{t}\equiv
A_0\pmod D$ (a discrete logarithm modulo $u$; a correct implementation must
enumerate the full preperiod-$s$ plus period $\mathrm{ord}_u(2)$ structure of
$2^{t}\bmod D$).
\end{lemma}
\begin{proof}
For $j<s$ every member has $\vtwo{\cdot}$ exactly $j$ and odd part
$q_0+2^{s-j}uk$, odd and affinely increasing $\ge q_0$; a power (odd part
$1$) occurs iff $q_0=1$.
\end{proof}

\subsection{L-exit roots: the class snapshot, the parity kill, and pairing}

\begin{proposition}[L-exit class-table snapshot; declared domain $10^{7}$
anchors]\label{prop:classes}
Within anchor indices $\le10^{7}$ (final value $n\le123{,}164{,}085\approx
1.23\times10^{8}$, $m=\vtwo{n+3}\le25$; collapsed-engine per-event walks,
cross-checked term-by-term against the uncollapsed step semantics on $600$
anchors, $0$ mismatches):
\textbf{(a)} there are $30{,}791{,}021$ root births, \textbf{none a power of
two}; composition (exact counts at $10^{6}$ anchors, independently reproduced
digit-for-digit): E1-reentries $2{,}787{,}614$ ($90.53\%$ \emph{of roots},
killed unconditionally by Lemma~\ref{lem:e1}; at the $10^{7}$ scale E1r is
covered by the lemma and the L-exits by the direct scan\footnote{The $10^{7}$
sweep (births7) performs the power check only on L-exit roots (the E1r branch
carries no check); the E1r-side evidence chain is Lemma~\ref{lem:e1}
(unconditional) plus the full $10^{6}$-anchor direct scan, $3\mid A$ in all
cases.}), L-exits $291{,}483$ ($9.47\%$; Lk $145{,}614$ $+$ L3 $145{,}869$).
\textbf{(b)} Every L-exit root falls into one of \textbf{20 classes} keyed by
the 4-segment key $(\mathrm{ty},\vtwo{A},A\bmod144,q\bmod144)$ ($=17$
3-segment coarse keys; the main chain L3$|$j2 splits into $4$ $q$-subclasses);
each class fits an arithmetic progression $\{A_0+Dk\}$ ($D$ = gcd of observed
differences; fit consistent with all observed members, independently
re-verified), with observed modulus family $D=2^{s}u$, $u\in\{9,45\}$,
$s\in\{4,7,11,12,13,16\}$. Representative classes ($A_0+Dk$): L3$|$j2
$=76+144k$; Lk$|$j0 $=77+144k$; Lk$|$j3 $=584+1152k$; Lk$|$j1
$=9230+18432k$; L3$|$j0 $=583+1152k$; L3$|$j4 $=36880+73728k$; Lk$|$j2
$=294932+589824k$; odd classes $A\equiv\{7,13,19,25,29,77,125\}\pmod{144}$;
late classes $A\equiv\{26,28,29\}\pmod{144}$ (first at anchors
$1{,}021{,}640$ / $2{,}043{,}278$ / $4{,}086{,}554$).
\textbf{(c) Pairing law.} Every Lk class $(\sigma,r_2)$ pairs with an L3
class $(\sigma{-}1,r_2)$ at the same $r_2$, with
$A_{\mathrm{Lk}}=A_{\mathrm{L3}}+1$ ($14{,}560/14{,}560$ over all $10^{5}$
anchors, zero violations; the converse fails for $7{,}305/14{,}587$ L3 roots,
which have no Lk partner).
\textbf{(d)} The table saturates after anchor $4{,}086{,}554$ (the $m{=}24$
record): $5.9\times10^{6}$ further anchors produce no new class. Cumulative
raw TM steps $\gtrsim10^{23}$ (re-estimated $\sim10^{24}$ from the verified
$16$--$22n^{2}$ cost band).
\end{proposition}

Certification layers (with the four corrections of the adversarial audit
birth-referee incorporated):
\begin{itemize}
\item \textbf{18 of the 20 classes are unconditionally safe by key
definition}: their residue $A\bmod144$ lies outside
$\{2^{t}\bmod144\}=\{1,2,4,8,16,32,64,128,112,80\}$, so \emph{no} member ---
present or future --- is a power of two.
\item \textbf{Lk$|$j3 ($A\equiv8\bmod144$) is a single-point collision}: its
only power, $8$, is below the least observed member $584$; and $A=8$ would
require $q\equiv1$, i.e.\ the \emph{unobserved} key (Lk,3,8,1), not this
class.
\item \textbf{L3$|$j4 ($A\equiv16\bmod144$) is the unique effective
collision class} (its power family $t\in\{4,10,16,\dots\}$ is unbounded): its
members form the family $A=9\cdot2^{12}W+16$ ($W$ odd $\ge1$,
$r_2-\sigma=30$) with odd part $2304W+1>1$ --- power-free \emph{conditional
on the family form} ($25/25$ orbital members at $10^{5}$ anchors $+$ $6/6$
synthetic members verified individually; direct scan $W\le2\times10^{6}$, $0$
powers).
\item \textbf{Artifact reality}: the offline certification artifact
(births11) contains \textbf{14} (3-segment-key) classes --- its member
sources omit the three late classes; those three are covered by mod-144
residue-class supersets (independent exhaustive brute-force decider).
\textbf{AP extrapolation direction:} when $D_{\mathrm{fit}}$ is a multiple of
the true period, the fitted progression is a \emph{subset} of the class
completion, not a superset --- a power-free fit does not certify a power-free
class (counterexample: the true class $\{80+48k\}$ contains $128$, while the
observations $80,176,272$ fit the power-free $\{80+96k\}$). The actual load
carried by the AP fit is consistency with all observed members (verified) and
the extrapolation of the L3$|$j4 family form (explicitly conditional).
\footnote{The power-containment decider powInAP of births11 shows
$1{,}722$ \emph{false negatives} in a $30{,}000$-case synthetic cross-check
(the $s{=}0$ branch is skipped entirely for odd $D$; the discrete-log loop
breaks out early) and $249$ false positives --- harmless for the $14$
certified classes (all have $j<s$, the defective branches untriggered;
independently re-verified by an enumeration-based decider) but \emph{not to
be reused as-is for new classes at $m>25$}; the correct tool is the
preperiod-plus-period enumerator of Lemma~\ref{lem:ap}.}
\end{itemize}

\begin{theorem}[parity kill of dangerous keys]\label{thm:parity}
Let $A=2^{j}q$ ($q$ odd) be any birth value; $A$ is a power of two iff $q=1$.
Hence: (i) $q\not\equiv1\pmod{144}\Rightarrow q\ne1\Rightarrow$ not a power;
(ii) $q\equiv1\pmod{144}\wedge q>1\Rightarrow q\ge145\Rightarrow$ not a power;
(iii) $j=0$ ($A$ odd) $\Rightarrow$ not a power. Counting with the full
4-segment keys, \textbf{19 of the 20 observed classes} have $q$-segment
$\not\equiv1\pmod{144}$ --- \emph{permanently safe by key definition}
(including Lk$|$j3$|$A8$|$q73); the unique $q\equiv1$ class is L3$|$j4 (its
family-level elimination is Proposition~\ref{prop:drift}). Every dangerous
key $(\cdot,j\ge1,2^j\bmod144,1)$ contains at most one power, $A=2^j$ (the
key fixes $\vtwo{A}=j$), excluded at the family level: members satisfy
$A=3\cdot2^{k-1}w+2^{j}\mu\ge3\cdot2^{k-1}\cdot3+2^{j}>2^{j}$ ($w\ge3$,
$\mu\ge1$).
\end{theorem}
\begin{proof}
A power of two has odd part $1$; (i)--(iii) are immediate, and the $19/20$
count is a key-segment check against Proposition~\ref{prop:classes}. The
magnitude inequality is as displayed.
\end{proof}

\begin{proposition}[m-record chain and class-generation timetable; declared
domain $10^{7}$ anchors]\label{prop:mrecords}
(Independent full re-run ref-b6, digit-for-digit.) The m-records $m=1..25$
all occur exactly at $n=3\cdot2^{m}-3$ (the pure $w=3$ doubling chain, direct
transfers $n\to2n{+}3$: $3\to9\to21\to\cdots\to100663293$), at anchor indices
$0,1,2,4,8,\dots,8{,}173{,}106$. \textbf{Of the first $25$ m-records, $14$
create a new class} ($m\in\{4,5,7,8,11,12,13,14,16,17,21,22,23,24\}$, plus
five additional instants at $w\in\{9,15,21\}$); the record anchors
$m=18,19,20,25$ create none. The table is unsaturated through $10^{7}$
anchors (last new class at anchor $4{,}086{,}554$); \emph{``the class table
grows without bound in $m$'' is a trend inference, not a theorem} --- the
per-$m$ visits of the $w=3$ chain guarantee unbounded \emph{opportunity} for
growth, while per-$m$ generation is content of Assumption~M.
\end{proposition}

\subsection{The drift law $\mathrm{M''}$: the final
explicit form of the residual}

Every L-exit root birth satisfies $A=\sigma+r_2/2+1=(3\sigma+d)/2+1$
($d=r_2-\sigma$). Synthetic experiments (constructed anchors $P(2^{a}w-3)$
walked directly in the collapsed engine, $w\equiv3\pmod6$) show that births
with $q\equiv1\pmod{144}$ occur only in \emph{clean geometric-head windows}
--- L-configuration run vectors
\begin{gather*}
\text{Lk-type }(\sigma=2^{k}w{+}1):\quad
[\,1,\ 2^{k}w{+}c,\ 2^{k}w,\ 2^{k+1}w,\ 2^{k+2}w{-}1\,],\\
\text{L3-relay }(\sigma=2^{k}w\ \text{exactly}):\quad
[\,1,\ 2^{k}w{+}c,\ 2^{k}w{-}1\,]\\
\text{or with geometric tail}\quad [\,1,\ 2^{k}w{+}c,\ 2^{k}w,\ 2^{k+1}w{-}1\,],
\end{gather*}
each window $k$ contributing an \emph{affine family}, linear in the window
parameter $w$:
\begin{equation}\label{eq:family}
A=3\cdot2^{k-1}\,w+2^{j}\mu,\qquad \mu\in\{1,145,289,\dots\}\ \text{odd
with}\ \mu\equiv1\!\pmod{144},\qquad d=2(m_0{-}1)-3\varepsilon,
\end{equation}
with $m_0=2^{j}\mu$ and $\varepsilon\in\{0,1\}$ the $+1$-offset of $\sigma$
(Lk windows $\varepsilon=1$, L3 windows $\varepsilon=0$; the values
$d=30,59,126,251,288,510,576,575$ all match). \emph{The coefficient
$2^{k-1}$ is a property of the window $k$, not of the class}: the class
L3$|$j4 in the form $9\cdot2^{12}W+16=3\cdot2^{12}w+16$ ($W=w/3$, odd
automatically since $w\equiv3\bmod6$) is specific to the $k=13$ window, and
its members occur only in anchor segments with $\vtwo{n+3}=13$. Write
\emph{$\mathrm{M'}$} (the residual arithmetic proposition) for: every L-exit
root birth $A=\sigma+r_2/2+1$ in every segment of the anchor sequence has
odd part $q\ge3$ (equivalently $A\notin\{2^{t}:t\ge2\}$).

\begin{proposition}[drift law; family census]\label{prop:drift}
(a) \emph{Verified family table} (census: $a=13..300\times w\in\{3,9,15\}$,
$864$ segments, all $1{,}683$ births with $q\equiv1$ \emph{fully classified},
each member verified against the exact identity; synthetic extension
$a=13..80\times6w$: $408$ segments, $7{,}062$ root births, all $336$
$q\equiv1$ births inside $\{F_{13},F_{26},F_{50}\}$, $0$ powers):
\begin{center}\footnotesize\setlength{\tabcolsep}{3.5pt}
\begin{tabular}{rllrlrll}
\toprule
window $k$ & type & key & $\sigma$-form & $d$ & $A$-form & first / persist. & members\\
\midrule
13 & L3 & (L3,4,16,1) & $2^{13}w$ & 30 & $3\cdot2^{12}w+2^4$ &
$a{=}13$ (exact) & 3\\
26 & Lk & (Lk,5,32,1) & $2^{26}w{+}1$ & 59 & $3\cdot2^{25}w+2^5$ &
$a\ge27$ (persist.) & 822\\
50 & L3 & (L3,6,64,1) & $2^{50}w$ & 126 & $3\cdot2^{49}w+2^6$ &
$a{=}50$ (exact) & 3\\
99 & Lk & (Lk,7,128,1) & $2^{99}w{+}1$ & 251 & $3\cdot2^{98}w+2^7$ &
$a\ge100$ (persist.) & 603\\
110 & L3 & (L3,0,1,1) odd & $2^{110}w$ & 288 & $3\cdot2^{109}w+145$ &
$a{=}110$ (exact) & 3\\
194 & L3 & (L3,8,256,1) & $2^{194}w$ & 510 & $3\cdot2^{193}w+2^8$ &
$a{=}194$ (exact) & 3\\
219 & L3 & (L3,0,1,1) odd & $2^{219}w$ & 576 & $3\cdot2^{218}w+289$ &
$a{=}219$ (exact) & 3\\
219 & Lk & (Lk,1,2,1) & $2^{219}w{+}1$ & 575 & $3\cdot2^{218}w+2\cdot145$ &
$a\ge220$ (persist.) & 243\\
\bottomrule
\end{tabular}
\end{center}
The odd part of each family equals $3\cdot2^{k-1-j}w+\mu$, odd: for $k{=}13$,
$\ge768w+1\ge2305$; for $k{}=26$, $\ge3\cdot2^{20}w+1\ge9{,}437{,}185$; the
same shape, larger, for $k{=}99/194$; the odd families ($j{=}0$) are
$\ge145/289$; the $k{=}219$-Lk family is $\ge3\cdot2^{217}w+145\ge145$ ---
\emph{all eight families are parity-killed} (all $1{,}683$ $q\equiv1$ members
pass).
\textbf{(b) Structural facts.} (1) The persistent families ($k=26,99,219$-Lk)
have $\sigma$ and $A$ independent of the anchor level $a$ (after first
appearance every segment hits the same $(\sigma,A)$, varying only with $w$)
--- mid-ladder windows (round-law locality: the $\alpha$-walk is autonomous
and triggers on passage); the level-exact families ($k=13,50,110,194,219$-L3)
occur only at levels whose endgame lands exactly in that window. (2) The
window positions $\{13,26,50,99,110,194,219\}$ (the eighth family first
appears at level $220$) form two near-doubling chains ($13\to26\to50\to99\to
194$ and $110\to219$) --- consistent with the timetable of the
$\vtwo{\cdot}$-record chaos (the non-periodicity clause of
Theorem~\ref{thm:noCycle}: new bits injected forever $\Rightarrow$ new
windows expected forever). (3) Beyond these eight families there is
\emph{no other} $q\equiv1\pmod{144}$ birth in the census or synthetic range;
$0$ powers; all E1r births divisible by $3$; all landings $\equiv3\pmod6$ and
$\ge n+6$. The constructed anchors $a\le80$ produce several further
$q\not\equiv1$ keys (e.g.\ Lk$|$j1$|$A62$|$q31, first at $a{=}49$) --- all
safe by key definition (Theorem~\ref{thm:parity}(i)).
\textbf{(c) $\mathrm{M''}$ (drift law).} For every $m\ge1$ and every
$w\equiv3\pmod6$, the $q\equiv1\pmod{144}$ L-exit root births in the segment
of the anchor $n=2^{m}w-3$ belong to finitely many window families
$\{F_k\}_{k\in K(m)}$ ($K$ monotone, $K\cap[1,300]=\{13,26,50,99,110,194,
219\}$), each of the form~\eqref{eq:family} with odd part odd and $\ge145$.
\textbf{(d) Equivalence statement.} $\mathrm{M''}\Rightarrow\mathrm{M'}$
(``no L-exit root birth on the orbit has odd part $q=1$''); conversely, given
the enumeration ring (every dangerous key belongs to some window family),
$\mathrm{M'}\Leftrightarrow\mathrm{M''}$. Relation to Assumption~M:
$\mathrm{M}\Rightarrow$ class preservation $\wedge\,\mathrm{M'}$
(Theorem~\ref{thm:main} $+$ the contrapositive of Theorem~\ref{thm:root}),
while $\mathrm{M'}\wedge$ class preservation $\Rightarrow$ non-halting
(Corollary~\ref{cor:assembly2}) --- so within the established architecture
$\mathrm{M''}$ (together with class preservation) \emph{carries exactly the
entire non-halting content of Assumption~M}: the endgame/round-law part of M
determines when windows activate, and $\mathrm{M''}$ compresses M into a
single arithmetic sentence in birth coordinates.
\end{proposition}

\begin{proposition}[synthetic assertion table for $m=26..40$]\label{prop:m2640}
Over $60$ segments ($15$ levels $\times\,w\in\{3,9,15,21\}$) of constructed
anchors walked directly, all assertions pass: $0$ power births ($60/60$);
all E1r divisible by $3$; all landings $\equiv3\pmod6$ with $\Delta\ge6$
($60/60$). At $a=26$: \emph{zero} $q\equiv1$ members (the $k{=}26$ window
not yet active; the $k{=}13$ key (L3,4,16,1) has no recurrence in $26..40$).
At $a=27..40$ ($56$ $q\equiv1$ members): \emph{56/56 satisfy the exact
$F_{26}$ identity} $A=3\cdot2^{25}w+32=9\cdot2^{25}W+32$ ($W=w/3\in
\{1,3,5,7\}$, all odd --- structural from $w\equiv3\pmod6$),
$\sigma=2^{26}w+1$, $r_2-\sigma=59$, $\vtwo{A}=5$, odd part
$q=3\cdot2^{20}w+1$ odd $\ge9{,}437{,}185>1$, hence not a power.
Cross-checks: the $a=13$ segment reproduces the orbital observations
term-by-term ($R(24576,24606)\to A=36880$; $R(12305,36886)\to A=30749$); the
orbital members $w=27,33,3501,6357$ also lie in $F_{13}$ (verified
digit-by-digit against $12288w+16$).
\end{proposition}

\begin{corollary}[conditional assembly; the $\mathrm{M''}$ route]\label{cor:assembly2}
Assume (i) $\mathrm{M'}$ (equivalently $\mathrm{M''}$ under the enumeration
ring) and (ii) class preservation (every anchor $n_k\equiv3\pmod 6$). Then
$M$, started on the blank tape, never halts.
\end{corollary}
\begin{proof}
By Theorem~\ref{thm:root} a halt requires a root birth equal to $2^{t}$.
Roots split into: E1-reentries (excluded unconditionally by
Lemma~\ref{lem:e1}); anchor initializations ($n$ odd $\Rightarrow$ first
sweep $y=n-2$ odd $\Rightarrow$ no birth --- guaranteed by class
preservation); L-exits (excluded by $\mathrm{M'}$).
\end{proof}
The class-preservation clause is \emph{necessary}: even-anchor births are an
independent genuine power channel --- at $n=10,42,170,682,\dots$ the anchor
birth $A=3n/2+1=16,64,256,1024,\dots$ is exactly a power of two, matching a
subfamily of the known even-$n$ halting set $\{10,26,42,\dots\}$.

\begin{remark}[the three precise gaps; unconditional resolution refused]\label{rem:gaps}
All assertion tables pass ($0$ powers, $56/56$, all eight families killed,
class tables and growth as stated), but the proof chain does not close ---
exactly three rings are missing:
\textbf{(1) unbounded window enumeration $=$ M:} the dangerous-key set grows
with $m$ (one family $\to$ eight by $a\le300$) with no evidence of stopping
(the $\vtwo{\cdot}$-record timetable $\Rightarrow$ new windows expected
forever); ``every family is parity-killed'' does not imply ``these are all
the families'' --- the $m\le2000$ table exhaustion plus synthetic $300$ is an
empirical closure.
\textbf{(2) class preservation is conditional} (L7/M-architecture,
Corollary~\ref{cor:assembly2}(ii)): even-anchor births are an independent
real power channel.
\textbf{(3) $m>25$ is a constructed-anchor property:} the synthetic
verification uses the collapsed engine (cross-validated by the births10
uncollapsed $0$-mismatch check $+$ the independent raw machine on $12/12$
anchors exactly), the constructed-anchor segments were not individually
re-checked at the raw level, and the ``orbital'' reading of $m=26..40$
depends on the orbit actually reaching those anchors (the m-record chain is
machine-confirmed through $m\le25$; beyond that belongs to M,
Proposition~\ref{prop:mrecords}).
\end{remark}

\begin{proposition}[depth-1 birth guards are not closed]\label{prop:depth1}
Let $I_1$ be $I_0$ together with same-shape algebraic guards at each birth
site (odd part of the birth value $\ge3$; the $f_2$ guard for $G(2,\cdot)$;
$w\equiv2\pmod3$ at $A=3$ sites; exact forms in the
\texttt{invariant-casecheck} suite). Then $I_1$ contains the entire real
orbit (trajectory library: $1{,}357{,}014$ states, zero violations), but
$I_1$ is \emph{not} closed: a closure search over the declared domain ($S/L$
tails of length $\le6$) yields $1{,}255$ violations, all of a single class
--- an odd-$r_2$ $L$ arming, via (L1), its S-odd child, whose local guard is
blind to the child's next-step birth.
\end{proposition}

\begin{proposition}[per-depth pollution of birth guards; witnessed]\label{prop:cascade}
For every tested depth $d\in\{1,2,3\}$ there is a configuration in $I_0$
whose pipeline has $d-1$ S-even births that are not powers of two, followed
by a power-of-two birth, and which truly halts:
\begin{center}
\begin{tabular}{clllr}
\toprule
depth & witness & birth sequence & phase events & raw steps\\
\midrule
$d{=}1$ & $L[1,3,1,3,1]$ & $8,4$ & 13 & 195\\
$d{=}1'$ & $S[1,1,1,1,1,1]$ & $8,4$ & 15 & 230\\
$d{=}1''$ & $G[4,2,7,5]$ & $16,8,4$ & 33 & 728\\
$d{=}2$ & $S[3,2,25,43,23,12,4,2]$ & $5,64$ & 152 & 10,313\\
$d{=}3$ & $S[12,28,23,42,14,22,38,31]$ & $27,71,128$ & 542 & 71,856\\
\bottomrule
\end{tabular}
\end{center}
All five witnesses were independently reproduced at the raw-machine level
(event-by-event tape agreement with the phase configurations, raw halt
confirmation, full chains inside $I_0$; additionally, the apparent
non-halting witness $S[39,34,5,30,43,19,37]$ of the $400$-step-capped search
halts within $442$ events under a larger cap --- a cap artifact, not a
counterexample). Moreover \emph{residue-class refinement is impossible}: the
bad births $\{2^{t}\}$ have residues mod $12$ in $\{4,8\}$, the very classes
occupied by the orbit's own S-even heads (the $A\bmod12$ histogram covers
$\{0,\dots,9\}$, with $9{,}282$ cases $A\equiv4$ and $396$ cases
$A\equiv8$); the intra-class split is exactly $\nu_2$ --- irreducibly
2-adic. The entry edges of bad births form a \emph{family of birth-entry
edges}: (L1) dominant (the $1{,}255$ violations of $I_1$ and the $d=2$
witness both enter through it), with (R) also occurring (the $d=3$ witness
enters through (R)$\to S_3$); the pollution mechanism is the per-level
translation of the (S1)/(L1) pipeline (Remark~\ref{rem:translate}).
\end{proposition}

\begin{remark}[structural translation for $d\ge4$; not proved]\label{rem:translate}
After each per-level repair (a depth-$d$ guard reads ``the first $d$ births
of the pipeline are non-powers-of-two''), an edge of the same shape
reappears one level deeper in the next $L$-round: each repair translates the
same-shaped pipeline by one level. Per-depth witnesses for every $d\ge4$ are
\emph{expected} from this mechanism (random search stably produces fresh
witnesses at each depth), but their per-$d$ existence is precisely the
statement at issue and is not claimed here; ``no finite guard depth closes''
is accordingly stated at witnessed level (Proposition~\ref{prop:cascade} plus
this mechanism), not as a theorem.
\end{remark}

\begin{remark}[relation to the no-precision conjecture]\label{rem:noprec2}
The companion derivation paper states the no-precision conjecture for coarse
templates (per-run residues at a fixed modulus plus linear inequalities);
Proposition~\ref{prop:cascade} does not prove it and does not change its
status. What it supports, at witnessed level (grade B), is the strengthening
direction beyond that language: residue classes and algebraic head-form
guards alike fail at every finite guard depth. The forcing edges are not on
the mod-3 side (Theorem~\ref{thm:mod3} seals that side completely) but at
the 2-adic birth sites: the entry edges form a family ((L1) dominant, (R)
also occurring), with the (S1)/(L1) pipeline translating one level deeper
per repair (Remark~\ref{rem:translate}).
\end{remark}

\section{Landmarks and classes}\label{s:landmark}

Fix an odd \emph{unit} $w\ge3$ for the rest of the paper.

\begin{definition}[landmark]\label{def:landmark}
For $\alpha\equiv1\pmod6$, $\alpha\ge7$, $t\ge9$, and slack $s\ge0$, let
$E:=2^{t-3}w$ and
\[
L(\alpha,t;s)\;=\;G\bigl(\alpha+3\cdot2^{t-3}w,\;\alpha-1+5\cdot2^{t-3}w,\;
2^{t}w,\;2^{t+1}w,\;\dots,\;2^{t+s}w,\;2^{t+s+1}w-1\bigr).
\]
Its \emph{class} is $h=\vtwo{\alpha+1}$.
\end{definition}

\begin{lemma}[class parametrization]\label{lem:class}
Fix $h\ge1$. The integers $\alpha\equiv1\pmod6$, $\alpha\ge7$ with
$\vtwo{\alpha+1}=h$ are exactly
\[
\alpha \;=\; 2^{h+1}(3c+\rho_h)+2^{h}-1,\qquad c\ge0,
\]
where $\rho_h\in\{0,1,2\}$ is the unique residue with
$(-1)^{h+1}\rho_h+(-1)^{h}\equiv1\pmod3$.
\end{lemma}

\begin{proof}
$\alpha+1=2^{h}\,\omega$ with $\omega$ odd forces
$\omega=2(3c+\rho)+1$ with $\rho\in\{0,1,2\}$, $c\ge0$, giving the shape.
Modulo $3$, using $2\equiv-1$: the displayed congruence is the condition
$\alpha\equiv1$, which singles out one $\rho$ since the three values of
$\rho$ give the three residues.
\end{proof}

\begin{theorem}[round law]\label{thm:round}
There exist functions $\Delta\alpha,\Delta t:\nn_{\ge1}\to\nn$ with
$\Delta\alpha(h)\equiv0\pmod6$ and $\Delta t(h)\in[2,9]$ such that for every
landmark $L(\alpha,t;s)$ of class $h$:
\begin{enumerate}
\item if $s\ge\Delta t(h)$, the phase trajectory returns to the landmark
$L\bigl(\alpha+\Delta\alpha(h),\,t+\Delta t(h);\,s-\Delta t(h)\bigr)$, and
no halting occurs;
\item if $s<\Delta t(h)$, the trajectory exits to the anchor
$P(n')$ with $n'=2^{t+s+1}w+c$, $c\equiv3\pmod6$, $c\ge3$.
\end{enumerate}
\end{theorem}

The proof runs through \S\ref{s:cascade}--\ref{s:master} by strong
induction on $h$; it is completed in \S\ref{s:master}.

\subsection{The halving cascade}\label{s:cascade}

\begin{lemma}[cascade]\label{lem:cascade}
Let $L(\alpha,t;s)$ have class $h$, write $\alpha+1=2^{h}\omega$ ($\omega$
odd), assume $t\ge h+4$, and put
\[
\Psi \;=\; \omega+3\cdot2^{t-3-h}w .
\]
Then $\Psi$ is odd, and the first $h$ events of the round are the halving
steps with leading heads
\[
a_i \;=\; 2^{h-i}\,\Psi-1,\qquad i=0,1,\dots,h,
\]
so that the $h$-th of them lands on the configuration
\[
S_0 \;=\; G\bigl(\Psi-1,\;\Psi-1,\;2\Psi,\;4\Psi,\;\dots,\;2^{h-1}\Psi,\;
\underbrace{\alpha+5E}_{\beta},\;2^{t}w,\;\dots,\;2^{t+s+1}w-1\bigr),
\]
i.e.\ on $\Phi(\Psi,h)$ with its final run $2^{h}\Psi-1$ replaced by the
triple $(\beta,\ \text{$E$-geometric tail},\ -1\ \text{terminator})$,
$\beta=\alpha+5E$.
\end{lemma}

\begin{proof}
The round opens with rule (G$\ge$) on the leading head
$a_0=\alpha+3\cdot2^{t-3}w$. Since $\alpha\equiv1\pmod6$ and
$3\cdot2^{t-3}w\equiv0\pmod6$ ($t\ge9$), $a_0$ is odd, and by the halving
arithmetic of Theorem~\ref{thm:prologue} the successive leading heads are
$a_i=(a_0-\delta_i)/2^{i}$-type where the affine bookkeeping gives
$a_i+1=2^{h-i}\Psi$: indeed $a_0+1=\alpha+1+3\cdot2^{t-3}w
=2^{h}\omega+3\cdot2^{t-3}w=2^{h}\Psi$ by definition of $\Psi$, and the
halving map divides $a+1$ by $2$ at each step while leaving the tail
invariant, provided each intermediate head is odd --- which holds because
$a_i+1=2^{h-i}\Psi$ with $\Psi$ odd. Meanwhile the displaced material
accumulates as the geometric tail $(2\Psi,4\Psi,\dots)$ exactly as in the
prologue. At $i=h$ the head is $\Psi-1$ and the configuration is $S_0$ as
stated; the run $\beta=\alpha+5E$ arises because the second leading run of
$L$ is $\alpha-1+5E$ and the cascade transfers the increment $+2$ net.
(Each such identity was verified symbolically and on
$h=1,\dots,60$ by exhaustive run-by-run comparison.)
\end{proof}

\begin{lemma}[entry to ladder scale]\label{lem:entryladder}
For $h\ge12$, the trajectory from $S_0$ reaches, after reading only runs up
to exponent $\le10<h-1$, the landmark-scale configuration
$L(7,9)$-type at unit $\Psi$ with the modified tail
$(\beta,E\text{-tail},-1)$ carried verbatim.
\end{lemma}

\begin{proof}
From $S_0$ the machine consumes the $\Psi$-geometric runs in the fixed
manner of a fresh $\Phi(\Psi,h)$-trajectory; the first landmark of any
$\Phi$-trajectory is $L(7,9)$-shaped (leading counter $7$, time parameter
$9$), reached having read only the runs $2\Psi,\dots,2^{9}\Psi$. The
modified tail sits entirely at exponents $\ge t>h+4\ge16>10$, hence beyond
the read window and carried verbatim. (Symbolic derivation; witnessed on
$h=12,\dots,60$.)
\end{proof}

\subsection{The universal ladder}\label{s:ladder}

\begin{definition}[ladder]\label{def:ladder}
Define the ladder as the sequence $(A_j,\tau_j)_{j\ge0}$ with
$(A_0,\tau_0)=(7,9)$ and
\[
(A_{j+1},\tau_{j+1}) \;=\; \bigl(A_j+\Delta\alpha(H_j),\;\tau_j+\Delta t(H_j)\bigr),
\qquad H_j:=\vtwo{A_j+1},
\]
where $\Delta\alpha,\Delta t$ are the functions of Theorem~\ref{thm:round}.
\end{definition}

\begin{lemma}[autonomy]\label{lem:autonomy}
On any $\Phi$-trajectory (with any admissible modified tail), the
successive landmarks follow the ladder: the $j$-th landmark is
$L(A_j,\tau_j;\cdot)$ at the unit scale, with the tail parameters evolving
independently. Consequently the sequence of classes visited depends only on
$j$, not on the unit or the tail.
\end{lemma}

\begin{proof}
By Theorem~\ref{thm:round} the passage landmark$\to$landmark changes
$(\alpha,t)$ by $(\Delta\alpha(h),\Delta t(h))$ with $h$ the class of the
current $\alpha$; the unit and tail enter only through the slack bookkeeping
(which decides landmark-vs-exit, not the map itself).
\end{proof}

\begin{lemma}[margins and class bound]\label{lem:ladder}
The ladder satisfies, for all computed positions:
\begin{enumerate}
\item $\tau_j\ \ge\ 2H_j+16$ for all $j\ge3$;
\item $A_j\ \le\ \tau_j^{2}$ for all $j\ge8$.
\end{enumerate}
Moreover, if these inequalities hold up to position $j$, they persist:
$A_{j+1}\le A_j+2H_j+6\le\tau_j^{2}+4\log_2\tau_j+6\le(\tau_j+2)^{2}\le\tau_{j+1}^{2}$
using $\Delta t\ge2$.
\end{lemma}

\begin{proof}
Induction as displayed, given the band bound
$\Delta\alpha(H)\le 2H+6$, which holds on all classes where the round law
has been established (this is part of the induction of \S\ref{s:master}:
the induction hypothesis at class $h$ includes the band bound for all
classes $\le h$). Items (1)--(2) were verified exhaustively on the first
$200{,}000$ ladder positions, which cover all $\tau_j\le10^{6}$.
\end{proof}

\begin{corollary}[well-foundedness]\label{cor:wellfounded}
For $h\ge16$, every ladder position with $\tau_j\le h-1$ satisfies
$H_j\le 2\log_2 h<h$: the ladder prefix traversed inside a round of class
$h$ involves only rounds of class strictly smaller than $h$.
\end{corollary}

\subsection{The master segment}\label{s:master}

Fix $h\ge16$ and a landmark $L(\alpha,t;s)$ of class $h$. By
Lemmas~\ref{lem:cascade}--\ref{lem:entryladder} the round consists of: the
cascade into $S_0$; a ladder prefix from $L(7,9)$-scale up to the position
$j_0$ = the last $j$ with $\tau_j\le h-65$ (all of whose rounds have class
$<h$ and are covered by the induction hypothesis); and the \emph{master
segment}: the events from the ladder literal
$L(A_{j_0},\tau_{j_0})$-scale to the end of the round.

\begin{assumption}[master closure]\label{ass:master}
For every $h\ge1$, the symbolic derivation from the ladder literal
$(A_{j_0(h)},\tau_{j_0(h)},h)$ with the outer parametrization of
Lemma~\ref{lem:class} ($c\ge0$, $u\ge0$, $w=6u+3$) and $t\ge2h+16$ is
\emph{single-path}: every branch decision (parity and run-size comparisons)
is forced uniformly over the entire domain, no split ever occurs, and the
derivation terminates in $O(\mathrm{poly}(\log h))$ events at either the
next real landmark (yielding $\Delta\alpha(h)\equiv0\bmod6$,
$\Delta t(h)\in[2,9]$) or, if $s<\Delta t(h)$, at the anchor exit of
Theorem~\ref{thm:round}(ii). Moreover the derivation is \emph{$h$-uniform}:
its complete form set depends on $h$ only through the ladder literals
$(A_{j_0},\tau_{j_0})$ and the exponent offsets of the modified tail
(the uniformity clause, spelled out in Remark~\ref{rem:uniform}); this is
what allows the finite audits below to be lifted from the verified
classes to all $h$.
\end{assumption}

\begin{remark}[the uniformity clause and its evidence]\label{rem:uniform}
The $h$-uniformity clause of Assumption~\ref{ass:master} states that the
master derivation depends on $h$ only through the ladder literals
$(A_{j_0},\tau_{j_0})$ and the exponent offsets of the modified tail,
whose margins grow monotonically with $h$ (the read window is bounded by
$\Delta t+3$ while the tail begins at exponent $t\ge2h+16$). Under this
clause, any failure of the single-path property at some $h_\ast>10^{6}$
would be witnessed by a violation of one of finitely many margin
inequalities, each monotone in $h$ beyond the verified range; the clause
was verified concretely on all classes $h\le2000$ (exhaustively) and on
sampled classes through $1.2\cdot10^{7}$. It is part of the assumption,
not a separate theorem: the monotonicity argument above is a structural
observation over the verified derivation shapes, and the safety proof
cites it as such (conditional on Assumption~\ref{ass:master}).
\end{remark}

\begin{proposition}[evidence for Assumption~\ref{ass:master}]\label{prop:master}
Assumption~\ref{ass:master} (single-path clause and uniformity clause
alike) was verified by exhaustive single-path symbolic
derivation: (i) for all $h=16,\dots,1024$, by direct agreement with the
independent per-$h$ direct engine ($1009/1009$); (ii) for $h=1025,\dots,2000$
exhaustively and for sampled $h$ up to $1.2\cdot10^{7}$ ($976$ table rows
plus on-the-fly ladder extensions, all single-path, all law-shape
consistent); (iii) at eight classes
$h\in[1028,1777]$ against numeric ground truth (exact; the class $h=1538$
among them realizes the rare value $\Delta t=9$, correcting the earlier
bound $[2,8]$ --- see the erratum in the dossier); (iv) at fresh
classes $h\in\{1823,2048,31337,666667\}$ --- at $h=1823$ with three-way
agreement (master $=$ direct engine $=$ precomputed table) --- and at
$h=45678$ via the published blind-test recipe, and --- with the ladder
recurrence computed on the fly rather than stored, cross-checked against
the stored ladder at $h=1777$ --- at the fresh classes
$h\in\{9999991,10^{7},10^{7}{+}3,12345679\}$, each single-path with
$\Delta\alpha/h\approx1.3382$.
\end{proposition}

\begin{proof}[Proof of Theorem~\ref{thm:round}]
By strong induction on $h$. \emph{Base $h\le15$:} direct single-path
symbolic derivations (the first $15$ rows of the per-$h$ table).
\emph{Step $h\ge16$:} the decomposition of \S\ref{s:cascade}--\ref{s:master}
expresses the round as cascade (Lemma~\ref{lem:cascade}, algebraic) $+$
entry (Lemma~\ref{lem:entryladder}) $+$ ladder prefix (rounds of class
$<h$, each correct by the induction hypothesis and
Corollary~\ref{cor:wellfounded}) $+$ master segment
(Assumption~\ref{ass:master}). The composite is single-path and ends in
outcomes (i) or (ii). The band bound $\Delta\alpha(h)\le2h+6$ required by
Lemma~\ref{lem:ladder} is carried along the same induction: it holds on the
base, and the master outputs satisfy it on all verified classes
($\Delta\alpha/h\approx1.338$ throughout, e.g.\ $\Delta\alpha(10^6)
=1338210$).
\end{proof}

\section{Endgames}\label{s:endgame}

\begin{theorem}[anchor exit]\label{thm:endgame}
Let $L(\alpha,t;s)$ of class $h$ occur on the $\Phi(w,m)$-trajectory
started from the anchor $P(n)$ ($n+3=2^{m}w$), and let $s<\Delta t(h)$.
Then $s$ is the slack to the \emph{terminator} run $2^{m}w-1$ of $\Phi$,
i.e.\ $t+s+1=m$, and the trajectory exits the $\Phi$-family at the anchor
\[
n' \;=\; 2^{m}\,w+c\;=\;(n+3)+c,\qquad c\equiv3\pmod6,\ c\ge3,
\]
so that $n'\ge n+6$.
\end{theorem}

\begin{proof}
Along a $\Phi$-trajectory the slack $s$ counts the geometric tail runs
beyond the current read window; the trajectory ends in an exit exactly
when the read window of the current round overruns the terminator
$2^{t+s+1}w-1$, which is the final run $2^{m}w-1$ of $\Phi(w,m)$; hence
$t+s+1=m$. The sweep arithmetic then resolves the configuration into a
two-run $R(n',n')$ with $n'=2^{m}w+c$. The congruence $c\equiv3\pmod6$
and $c\ge3$ were verified symbolically for every endgame cell (all
$h\le160$ exhaustively --- $841/841$ cells, including two independent
re-audits; the full $h\le1024$ grid adds $5049$ further cells, all passing,
for the ledger total $5890$), with the endgame anchor form additionally
matched against constructed landmark walks ($216/216$). The growth inequality follows:
$n'-(n+3)=c\ge3$, i.e.\ $n'\ge n+6$.
\end{proof}

\section{No halting}\label{s:safety}

\begin{theorem}[safety]\label{thm:safety}
No configuration on any trajectory from $P(3)$ is of the form
$G(1,\cdot)$ or $G(3,\cdot)$; in particular no $G$-state has leading head
$a\in\{1,3\}$, and by Lemma~\ref{lem:halt} $M$ never enters the halting
transition. (The master-segment part of the computational audit below is
conditional on Assumption~\ref{ass:master} in the standard way.)
\end{theorem}

\begin{lemma}[exhaustiveness of the state classification]\label{lem:exhaust}
Every $G$-state occurring on a trajectory from $P(3)$ occurs inside one of
the segment derivations of Theorem~\ref{thm:prologue},
Lemmas~\ref{lem:cascade}--\ref{lem:entryladder}, the rounds of
Theorem~\ref{thm:round} (base table for $h\le15$; ladder prefix rounds and
master segment for $h\ge16$), or the endgame cells of
Theorem~\ref{thm:endgame}; and each such segment derivation is single-path,
hence enumerates its events exactly.
\end{lemma}

\begin{proof}
By the assembly recursion of \S\ref{s:assembly}, the trajectory from $P(3)$
is the concatenation of the segment derivations of the anchor transitions.
Each segment is single-path: the prologue and cascade are deterministic
algebra (Theorem~\ref{thm:prologue}, Lemma~\ref{lem:cascade}); the base
rounds $h\le15$, the ladder-prefix rounds (classes $<h$, induction), the
master segments (Assumption~\ref{ass:master}), and the endgame cells are
single-path symbolic derivations by construction/verification. A
single-path derivation from a fixed initial configuration enumerates
exactly the events that occur.
\end{proof}

\begin{proof}[Proof of Theorem~\ref{thm:safety}]
We split the argument into an algebraic part and a computational part.

\emph{Algebraic part (prologue, entry, and cascade families).}
Let $a$ be the leading head of a $G$-state in:
(a) the prologue halving chain, where $a_i+1=2^{m-i}w$ with $w\ge3$ odd
(note $w\ge3$ since $n\equiv3\pmod6$ forces $n+3\ge6$); the prologue's
initial head $n+2\ge5$ is odd and $\ne3$ as $n\ge3$;
(b) the halving cascade of a round with class $h$ and $t\ge h+4$, where
$a_i+1=2^{h-i}\Psi$ with $\Psi=\omega+3\cdot2^{t-3-h}w$;
here $\omega\ge1$ is odd and $3\cdot2^{t-3-h}w\ge3\cdot2\cdot3=18$, so
$\Psi\ge19$ is odd;
(c) a landmark $G$-head $\alpha+3\cdot2^{t-3}w$ with $\alpha\equiv1\pmod6$,
$\alpha\ge7$, and $3\cdot2^{t-3}w\ge3\cdot64\cdot3>12$: the head is odd
(odd $+$ even) and $\ge19$.
In cases (a)--(c), if $a$ is odd then $a+1=2^{j}\cdot(\text{odd}\ge3)$
for some $j\ge1$, while $a\in\{1,3\}$ would require $a+1\in\{2,4\}
=2^{j}\cdot1$ with odd part $1<3$: contradiction. The even values that
occur ($a=w-1$ with $w\ge3$ odd) are $\ge2$ and even, hence also not in
$\{1,3\}$. So $a\notin\{1,3\}$ throughout these families.

\emph{Computational part (intermediate and master/endgame states).}
The $G$-states inside sweeps, master segments, and endgames are exactly
those enumerated by the single-path derivations of
Lemma~\ref{lem:exhaust}. The $2$-adic audit (exact two-variable $(u,c)$
solvers) checked, over $273{,}654$ symbolic forms --- all round classes
$h\le1024$, all endgame cells, and the full master-path form set (which,
by the uniformity clause of Assumption~\ref{ass:master}
(Remark~\ref{rem:uniform}), is $h$-uniform) --- that no form has a
leading head with $a+1$ a power of $2$ (this excludes $a\in\{1,3,7,15,
\dots\}$ simultaneously, in particular both halting heads) or reaches
such a head under any admissible continuation of its symbolic domain:
$0$ dangers, $0$ undecided.

By Lemma~\ref{lem:exhaust} the two parts cover every $G$-state on the
trajectory, and in neither part does a head in $\{1,3\}$ occur. By
Lemma~\ref{lem:halt}, $M$ never enters the halting transition.
\end{proof}

\section{Assembly and proof of the main theorem}\label{s:assembly}

\begin{proof}[Proof of Theorem~\ref{thm:main}]
By Lemma~\ref{lem:entry} the run reaches $P(3)$ at step $54$. Define the
anchor sequence $n_0=3$ and, inductively, $n_{k+1}:=n'$ the exit anchor of
Theorem~\ref{thm:endgame} applied to the $\Phi(w_k,m_k)$-trajectory started
from $P(n_k)$ (which exists by Theorem~\ref{thm:prologue} with
$n_k+3=2^{m_k}w_k$). We maintain the invariant $6\mid n_k+3$. It holds
for $n_0=3$. Suppose $6\mid n_k+3=2^{m_k}w_k$; then $3\mid w_k$, and
$6\mid 2^{m_k}w_k$ (both $2$ and $3$ divide it). The exit of
Theorem~\ref{thm:endgame} gives
\[
n_{k+1}+3 \;=\; 2^{m_k}w_k+(c_k+3),\qquad 6\mid c_k+3\ \ (\text{since } c_k\equiv3\pmod6),
\]
so $6\mid n_{k+1}+3$. In particular every $n_k\equiv3\pmod6$ (odd), which
is the domain of Theorem~\ref{thm:prologue}.
Theorem~\ref{thm:round} gives the complete decomposition of every
transition into prologue $+$ rounds $+$ exit, with
$n_{k+1}\ge n_k+6$ (Theorem~\ref{thm:endgame}); the sequence is defined for
all $k$; and by Theorem~\ref{thm:safety} no halting configuration ever
occurs. Since the phase calculus of Lemma~\ref{lem:rules} mirrors the
machine exactly (including the halting transition), $M$ never halts.
\end{proof}

\section{Computational components}\label{s:comp}

The following components are carried by exhaustive symbolic computation and
are stated precisely so that they can be re-audited or formalized
independently (scripts in the repository; each row is a one-line rerun):

\begin{center}\small
\begin{tabular}{@{}p{5.0cm}p{3.9cm}p{6.6cm}@{}}
\toprule
Component & Domain & Result\\
\midrule
Rule-level grid (Lemma~\ref{lem:rules}) & $a\le20\times b\le10\times$ tails & $564/564$ (\texttt{rulecheck})\\
Trajectory faithfulness (Lemma~\ref{lem:rules}) & $956$ anchors ($n{=}2..957$) & $955$ equal; $R(2,2)$ off-grammar\\
Prologue (Theorem~\ref{thm:prologue}) & $n\le6003$; $m\le1200$ & $1052/1052$\\
Cascade/entry (Lemmas~\ref{lem:cascade},\ref{lem:entryladder}) & $h\le60$ & $327/327$\\
Ladder margins (Lemma~\ref{lem:ladder}) & $j<200000$ & $199997{+}199992$\\
Base round law & $h\le15$ & single-path\\
Induction step & $h=16..1024$ & $1009/1009$\\
Master extension & $h=1025..1.2\cdot10^{7}$ & $976$ rows $+$ 4 fresh\\
Endgame cells (Theorem~\ref{thm:endgame}) & $h\le160$ ($\le1024$) & $841/841$ ($5890$ total)\\
Endgame forms vs.\ walks & $216$ cells & $216/216$\\
$2$-adic safety audit & $273654$ forms & $0$ dangers\\
Ledger closure (Theorem~\ref{thm:mod3}) & boundary $+$ $2^{40}$ random & $219{,}398$ author $+$ $20{,}508$ independent steps ($k\le10$), $0$ violations\\
Ledger orbit (Theorems~\ref{thm:mod3}, \ref{thm:haltpow}) & $10^{7}$ events & $0$ violations; $2{,}908$ births, $0$ powers\\
Independent raw re-check & $5\cdot10^{5}$ events & $647$ births, $0$ powers\\
Root partition (Theorem~\ref{thm:root}; ref-b1) & $2\times10^{7}$ uncollapsed events, $572$ anchors & $0$ provenance/inheritance/odd-part violations; raw $12/12$ anchors\\
E1-reentry scan (Lemma~\ref{lem:e1}; ref-b2) & $10^{6}$ anchors & $2{,}787{,}614/2{,}787{,}614$ with $3\mid A$; $3{,}669{,}365$ re-entries all $b\equiv2\bmod3$\\
Birth census (Proposition~\ref{prop:classes}; births5--7) & $10^{7}$ anchors & $30{,}791{,}021$ roots, $0$ powers (L-exit direct; E1r by lemma $+$ $10^{6}$ scan)\\
Uncollapsed cross-check (births10) & $600$ anchors, $9.5\times10^{6}$ events & $4{,}322$ births, $0$ mismatches\\
Class certification (Proposition~\ref{prop:classes}; births11, ref-b3/b4) & $20$ classes & $14$ AP-certified $+$ $3$ residue supersets; Lk$|$j3 single point; L3$|$j4 family ($W\le2\times10^{6}$, $0$ powers); pairing $14{,}560/14{,}560$\\
m-records (Proposition~\ref{prop:mrecords}; ref-b6) & $10^{7}$ anchors & $m=1..25$ all at $3\cdot2^{m}-3$; $14/25$ create classes\\
Family census (Proposition~\ref{prop:drift}; v7-census) & $a=13..300\times\{3,9,15\}$ & $864$ walks; $1{,}683$ $q\equiv1$ all in $8$ families; parity kill all\\
Synthetic families (v7-synth) & $a=13..80\times6w$ & $408$ walks; $7{,}062$ births; $336$ $q\equiv1$ in $\{F_{13},F_{26},F_{50}\}$; $0$ powers\\
$m{=}26..40$ assertions (Proposition~\ref{prop:m2640}) & $15$ levels $\times4w$ & $60$ segments: $0$ powers; $56/56$ exact $F_{26}$ form\\
Depth-1 guard search (Proposition~\ref{prop:depth1}) & $S/L$ tails $\le6$ & $1255$ violations, one class\\
Cascade witnesses (Proposition~\ref{prop:cascade}) & $5$ chains & $13/15/33/152/542$ events, raw halt\\
\bottomrule
\end{tabular}
\end{center}

\paragraph{The residual in its final form.}
Magnitude decisions are settled unconditionally by dominance ($h\ge82$;
$h\le81$ pointwise); $\Delta t(h)\in[2,9]$ on all verified classes
(Theorem~\ref{thm:round}); the halting branch is
unreachable and perturbation-robust ($2980/2980$ single/sparse-flip
counterfactuals close safely; adversarial profiles diverge but are
$2$-adically unrealizable and never halt). What remains is
Assumption~\ref{ass:master} beyond the exhaustive range ($h\le2000$
exhaustive, sampled to $1.2\cdot10^{7}$). A finite-state reduction
(\texttt{sigdyn.mjs}, \texttt{brent32.mjs}) shows this residual is exactly
the continuation of the \emph{record timetable} of the autonomous walk:
first arrivals of new class depths occur at rounds climbing by a factor
$\approx4$ per level (measured schedule $17@19{,}024$, $19@76{,}096$,
$21@304{,}383$, $22@1{,}217{,}531$, $24@4{,}870{,}122$, and depth $\ge32$
by step $1{,}246{,}751{,}041$ of the residue kernel, which terminates at a
bad residue --- the record-$32$ event --- with no cycle). In birth
coordinates the same residual takes the explicit form
$\mathrm{M''}$ (Proposition~\ref{prop:drift}): the dangerous-key families
grow with the window index exactly as the record timetable prescribes, and
``never halts $\Leftarrow$ $\mathrm{M''}$ $+$ class preservation''
(Corollary~\ref{cor:assembly2}) matches the master-closure chain in
non-halting content (Remark~\ref{rem:gaps} for the precise three-ring
difference).

\begin{theorem}[periodicity excluded on the verified slice; conjecture beyond]\label{thm:noCycle}
On the $32$-bit slice of the signature dynamics (the residual kernel of
the $K$-slice model at $K=32$): Brent cycle detection over
$1.25\cdot10^{9}$ steps finds no cycle, and the kernel terminates at a bad
residue (the record-$32$ event); the ladder residues over
$2.4\cdot10^{6}$ positions are pairwise distinct. Within this
architecture, any $K$-bit cycle of the kernel would suppress the record
timetable below depth $K$ forever, contradicting the observed climb on
the verified slice. Consequently no finite closure of
Assumption~\ref{ass:master} exists at $K\le32$; the corresponding
statement for $K>32$ --- equivalently, that the full signature dynamics
is not eventually periodic --- is a \emph{conjecture} (the record climb
is measured, not proved, beyond the slice). The residual is equivalent to
the $\nu_2$-chaos of the autonomous walk --- the irreducible core the
architecture avoids needing everywhere except at record events.
\end{theorem}

Formalizing Lemma~\ref{lem:rules} and the dominance inequality in a proof
assistant would raise the conditional theorem and the envelope to
proof-assistant grade; eliminating Assumption~\ref{ass:master} itself
requires, by Theorem~\ref{thm:noCycle}, mathematics beyond this
architecture (or unbounded per-class computation).

\end{document}
